\documentclass{article}
\usepackage[utf8]{inputenc}
\usepackage{amsmath}
\usepackage{amssymb}
\usepackage{geometry}
\geometry{a4paper, margin=1in}
\title{Reconstruction of Neural Tangent Kernel (NTK) Calculations for an MMNN Network}
\author{Analysis of provided notes}
\date{\today}

\begin{document}

\maketitle

\section{Introduction and Definitions}
This document first introduces the Multi-component and Multi-layer Neural Network (MMNN) architecture in its general n-dimensional form. It then reconstructs a series of calculations aimed at deriving the neural tangent kernel (NTK) for this network. The central hypothesis for the NTK calculation, in accordance with learning strategy S1 from the source paper, is that only the parameters of the external linear combination layer are trained, while the weights and biases internal to the activation function are fixed and randomly initialized.

\subsection{MMNN Architecture}
A single layer of an MMNN is a shallow neural network approximating a vector-valued function $\bmh:\mathbb{R}^{d_{in}} \to \mathbb{R}^{d_{out}}$. It is defined as:
\begin{equation}
    \bmh(\bm{x}) = \bm{A} \sigma(\bm{W}\bm{x} + \bm{b}) + \bm{c}
\end{equation}
This can be written in a more compact form to better illustrate its structure. By defining augmented matrices and vectors that encapsulate the biases, we have:
\begin{equation}
    \bmh(\bm{x}) = \tilde{\bm{A}}
    \begin{bmatrix}
        \sigma(\tilde{\bm{W}}\tilde{\bm{x}}) \\
        1
    \end{bmatrix}
\end{equation}
where the augmented components are defined as:
\begin{itemize}
    \item $\tilde{\bm{x}} = [\bm{x}^{\top}, 1]^{\top}$: The input vector augmented with a 1 to handle the bias term.
    \item $\tilde{\bm{W}} = [\bm{W}, \bm{b}]$: The weight matrix augmented with the bias vector $\bm{b}$, such that the product $\tilde{\bm{W}}\tilde{\bm{x}}$ correctly computes $\bm{W}\bm{x} + \bm{b}$.
    \item $\tilde{\bm{A}} = [\bm{A}, \bm{c}]$: The output matrix augmented with the output bias vector $\bm{c}$.
\end{itemize}

We call each element of $\bmh$, i.e., $\bmh[i] = \tilde{\bm{A}}[i, :] \cdot ([\sigma(\tilde{\bm{W}}\tilde{\bm{x}})]^{\top}, 1)^{\top}$, a \textit{component}. This representation highlights several key features:
\begin{enumerate}
    \item Each component is a linear combination of a shared set of basis functions, $\{\sigma((\bm{W}\bm{x}+\bm{b})_j)\}_{j=1}^n$.
    \item Different components of $\bmh$ use the same basis functions but with different linear coefficients, given by the rows of $\tilde{\bm{A}}$.
    \item According to learning strategy S1, only the external coefficients $\tilde{\bm{A}}$ are trained, while the internal parameters $\tilde{\bm{W}}$ that define the basis are randomly initialized and fixed.
    \item The output dimension $d_{out}$ (the number of components, or rank) and the network width $n$ (the number of basis functions) are tunable hyperparameters. Typically, $d_{out} \ll n$.
\end{enumerate}

In contrast to a typical FCNN layer, which takes the form $\sigma(\tilde{\bm{W}}\tilde{\bm{x}})$ and in which all weights are trained, an MMNN layer is composed of multiple components where each is a linear combination of randomly parameterized hidden neurons. This structure acts as a smooth decomposition that can be trained more effectively.

The decomposition theorems presented earlier provide a mathematical intuition for this architecture. The fixed, random basis functions generated by $\sigma(\tilde{\bm{W}}\tilde{\bm{x}})$ play a role analogous to the piecewise linear transformations $\psi_i$. The trainable parameters $\tilde{\bm{A}}$ then learn the appropriate linear combination of these basis functions, conceptually similar to learning the smoother functions $f_i$ in the decomposition. This parallel between the parameters $(\tilde{\bm{A}}, \tilde{\bm{W}})$ and the decomposition functions $(f_i, \psi_i)$ will be explored further in the report.

An MMNN is a full multi-layer composition of such blocks:
\begin{equation}
    \bmh_{total} = \bmh_m \circ \cdots \circ \bmh_2 \circ \bmh_1
\end{equation}
where each block $\bmh_i: \mathbb{R}^{d_{i-1}} \to \mathbb{R}^{d_i}$ is a multi-component shallow network. The overall network is characterized by its width (largest number of basis functions), rank (largest intermediate dimension), and depth (number of layers).

\subsection{NTK Preliminaries}
For the NTK calculation, we consider the trainable parameters $\theta = \{A, c\}$. The kernel is defined by the dot product of the network output gradients with respect to these parameters:
\begin{equation}
    K_{NTK}(x, x') = \mathbb{E}_{\theta_{init}} \left[ \left\langle \nabla_{\theta} h(x), \nabla_{\theta} h(x') \right\rangle \right]
\end{equation}
For a single layer, the gradient with respect to the coefficients of $A$ is the vector of activations $\sigma(Wx+b)$, and the gradient with respect to $c$ is a vector of 1s.

\section{Multi-component and Multi-layer Decomposition}
The effectiveness of MMNNs stems from their ability to decompose a complex function into a combination of smoother, more manageable components. This "divide and conquer" strategy is formalized below, first in one dimension and then generalized to higher dimensions.

\subsection{One-Dimensional Decomposition}
Given a function $f$ on an interval $[x_0, x_n]$ partitioned by points $x_0 < x_1 < \cdots < x_n$, we can decompose $f$ using piecewise linear transformations. Let $\calL_i:[a_i,b_i]\to [x_{i-1},x_i]$ be a linear map. We define a smoother version of $f$ on each subinterval as $f_i = f \circ \calL_i$.

The transformation $\psi_i(x)$ maps the subinterval $[x_{i-1}, x_i]$ to $[a_i, b_i]$ and is constant outside. It can be constructed with ReLU functions:
\begin{equation}
    \psi_i(x)=s_i\cdot \text{ReLU}\left( {x-x_{i-1}} \right)
    -s_i\cdot \text{ReLU}\left( {x-x_{i}} \right)+a_i.
\end{equation}
where $s_i = (b_i-a_i)/(x_i-x_{i-1})$.

\begin{theorem}[1D Decomposition]
The target function $f:[x_0,x_n]\to\mathbb{R}$ can be decomposed as:
\begin{equation}
    f(x) = \sum_{i=1}^n f_i\circ \psi_i(x) - \sum_{i=1}^{n-1}f(x_i)
\end{equation}
for any $x \in [x_0, x_n]$.
\end{theorem}
This theorem shows that a complex function can be represented by a sum of compositions of smoother functions ($f_i$) and simple piecewise linear maps ($\psi_i$), which are easily representable by a shallow network.

\subsection{N-Dimensional Decomposition}
The decomposition strategy extends to multiple dimensions by applying the 1D decomposition iteratively for each dimension. For a 2D function $f(x,y)$ on a domain $[x_0,x_n]\times [y_0,y_m]$, the decomposition formula is derived from the principle of inclusion-exclusion (also known as Poincaré's formula), which provides a way to express the function in terms of its values on sub-domains and at boundary points. This leads to the following theorem.

\begin{theorem}[2D Decomposition]
A function $f:[x_0,x_n]\times [y_0, y_m]\to\mathbb{R}$ can be expressed as:
\begin{equation}
    \begin{split}          
        f( x, y) & = \sum_{i=1}^n \sum_{j=1}^m f_{i,j}\Big( \psi_i(x), \phi_j(y)\Big) - \sum_{i=1}^n\sum_{j=1}^{m-1} f_{i,0}\Big(\psi_i(x), y_{j}\Big) \\
        & \quad - \sum_{i=1}^{n-1} \sum_{j=1}^m f_{0,j}\Big( x_i, \phi_j(y)\Big) + \sum_{i=1}^{n-1}\sum_{j=1}^{m-1} f (x_i, y_{j})
    \end{split}
\end{equation}
where $f_{i,j}$, $f_{i,0}$, $f_{0,j}$ are rescaled versions of $f$, and $\psi_i, \phi_j$ are the 1D piecewise linear transformations for each dimension.
\end{theorem}
This formula, based on the inclusion-exclusion principle, generalizes to any number of dimensions, as detailed in the following theorem. The proof of this general theorem can be established by recursion on the dimension $d$.

\begin{theorem}[N-Dimensional Decomposition]
The decomposition for a function $f: \prod_{k=1}^d [x_{k,0}, x_{k, n_k}] \to \mathbb{R}$ on a hyperrectangle is given by:
\begin{equation}
    f(\bm{x}) = \sum_{S \subseteq \{1, \dots, d\}} (-1)^{d-|S|} \sum_{\bm{j}_{\bar{S}}} \left( \mathcal{D}_S \left[ f(\cdot, \bm{x}_{\bar{S}, \bm{j}_{\bar{S}}}) \right] (\bm{x}_S) \right)
\end{equation}
where:
\begin{itemize}
    \item $\bm{x}_S = (x_k)_{k \in S}$ are the variables for dimensions in the subset $S \subseteq \{1, \dots, d\}$.
    \item $\bm{x}_{\bar{S}, \bm{j}_{\bar{S}}} = (x_{k, j_k})_{k \in \bar{S}}$ are grid points for dimensions in the complement $\bar{S}$.
    \item $\mathcal{D}_S$ is the decomposition operator applied to variables in $S$, which rescales the function on each subdomain of the partition, i.e., $\mathcal{D}_S[g](\bm{y}_S) = \sum_{\bm{i}_S} g_{\bm{i}_S}((\psi_{k, i_k}(y_k))_{k \in S})$.
\end{itemize}
This provides a mathematical foundation for the ability of MMNNs to handle high-dimensional, complex functions by breaking them into simpler parts.
\end{theorem}

\subsection{Further Considerations on Scope}
While the MMNN architecture can be extended to include residual connections (ResMMNNs), the corresponding NTK calculations become significantly more complex. Therefore, this analysis will focus on the standard MMNN architecture, with ResMMNNs being a subject for future work.

Furthermore, our analysis relies on a specific parameter initialization strategy tailored to the MMNN structure and the infinite-width regime ($n \to \infty$):
\begin{itemize}
    \item The inner weights $\tilde{\bm{W}}$, which define the basis functions, are drawn from a fixed, order-one distribution (i.e., their variance does not scale with $n$). This is a crucial and non-standard assumption. The rationale is that these weights define a partition of the input space. To achieve a progressively finer discretization of the domain as the width $n$ grows, the basis functions must be distributed across the entire domain, which requires their parameters to be of order one.
    \item The outer, trainable weights $\tilde{\bm{A}}$, which determine the coefficients of the basis functions, are scaled by $1/\sqrt{n}$ (i.e., initialized with a variance of $1/n$). This normalization is essential for the NTK theory: it ensures that the network output, being a sum over $n$ components, converges to a well-behaved Gaussian process in the infinite-width limit. This allows for a stable approximation of the function's values on the discretized domain created by $\tilde{\bm{W}}$.
\end{itemize}
To begin the analysis, we will focus on the Rectified Linear Unit (ReLU) as the activation function. Its property of being positive-homogeneous (i.e., $\text{ReLU}(\alpha x) = \alpha \text{ReLU}(x)$ for $\alpha > 0$) significantly simplifies the kernel computations, making it an ideal choice for the initial derivation.

\section{NTK Computation for a Single-Layer MMFN}
We now derive the Neural Tangent Kernel for a single-layer MMFN. For simplicity, we consider a single scalar output, such that the network is defined as:
\begin{equation}
    h(\bm{x}) = \frac{1}{\sqrt{n}} \sum_{j=1}^n A_j \sigma(\bm{w}_j^\top \bm{x} + b_j) + c
\end{equation}
The trainable parameters are $\theta = \{A_j\}_{j=1}^n \cup \{c\}$. The parameters $\{\bm{w}_j, b_j\}$ are fixed and randomly drawn, with $\bm{w}_j \sim \mathcal{N}(0, \Sigma_w)$ and $b_j \sim \mathcal{N}(\mu_b, \sigma_b^2)$.

\subsection{From Gradient Products to an Expectation}
In the infinite-width limit ($n \to \infty$), the NTK converges to an expectation over the distribution of the random parameters $(\bm{w}, b)$:
\begin{equation}
\label{eq:ntk_expectation}
    K_{NTK}(\bm{x}, \bm{x}') = \mathbb{E}_{\bm{w},b} \left[ \sigma(\bm{w}^\top \bm{x} + b)\sigma(\bm{w}^\top \bm{x}' + b) \right] + 1
\end{equation}

\subsection{Integral Formulation and Parameter Mapping}
The expectation in \eqref{eq:ntk_expectation} is explicitly an integral over the distributions of $\bm{w}$ and $b$:
\begin{equation}
    \mathbb{E}_{\bm{w},b}[\cdot] = \int_{\mathbb{R}^d} \int_{\mathbb{R}} \sigma(\bm{w}^\top \bm{x} + b)\sigma(\bm{w}^\top \bm{x}' + b) p(\bm{w})p(b) \,db \,d\bm{w}
\end{equation}
where $p(\bm{w})$ and $p(b)$ are the Gaussian probability density functions. A change of variables is performed from the parameter space $(\bm{w}, b)$ to the pre-activation space $(g_1, g_2)$, where $g_1 = \bm{w}^\top \bm{x} + b$ and $g_2 = \bm{w}^\top \bm{x}' + b$. This transforms the high-dimensional integral into an expectation over a 2D bivariate Gaussian distribution. The parameters of this distribution are:
\begin{itemize}
    \item \textbf{Means:} $\mu_1 = \mathbb{E}[g_1] = \mathbb{E}[\bm{w}^\top\bm{x} + b] = \mu_b$. By symmetry, $\mu_2 = \mu_b$.
    \item \textbf{Variances:} $\sigma_1^2 = \text{Var}(g_1) = \text{Var}(\bm{w}^\top\bm{x}) + \text{Var}(b) = \bm{x}^\top \Sigma_w \bm{x} + \sigma_b^2$. Similarly, $\sigma_2^2 = {\bm{x}'}^\top \Sigma_w \bm{x}' + \sigma_b^2$.
    \item \textbf{Covariance:} $\text{Cov}(g_1, g_2) = \mathbb{E}[(g_1-\mu_b)(g_2-\mu_b)] = \mathbb{E}[(\bm{w}^\top\bm{x} + b-\mu_b)(\bm{w}^\top\bm{x}' + b-\mu_b)] = \bm{x}^\top \Sigma_w \bm{x}' + \sigma_b^2$.
\end{itemize}
The problem is thus reduced to calculating the expectation of a product of ReLUs applied to non-centered, correlated Gaussian variables.

\subsection{Comparison with Standard 1-Layer FCNN Kernel}
The formulation of the MMFN kernel in \eqref{eq:ntk_expectation} differs crucially from the standard NTK derived for a simple 1-layer FCNN. In many standard FCNN analyses, the bias term is either omitted from the ReLU for simplicity or treated as a separate, trainable parameter whose contribution to the kernel is a simple additive constant. Consequently, the core expectation is often $\mathbb{E}_{\bm{w}}[\sigma(\bm{w}^\top\bm{x})\sigma(\bm{w}^\top\bm{x}')]$, which involves centered Gaussian variables.

For the MMFN, the bias $b$ is a random, non-trainable parameter *inside* the activation function. Its presence makes the pre-activations $g_1, g_2$ non-centered (if $\mu_b \neq 0$), leading to the more complex non-centered expectation derived below. This is a defining feature that captures the richer representation allowed by the random bias.

\subsection{Final Analytical Expression}
The NTK for a single-layer MMFN is given by the expectation over the product of two non-centered, correlated ReLU activations. This has a closed-form analytical solution, derived by first considering the joint Gaussian distribution of the pre-activations $g_1 = \bm{w}^\top \bm{x} + b$ and $g_2 = \bm{w}^\top \bm{x}' + b$.

\begin{theorem}[NTK for a Single-Layer ReLU MMFN]
Let the pre-activations $(g_1, g_2)$ be jointly Gaussian with parameters $(\mu_1, \sigma_1^2)$ and $(\mu_2, \sigma_2^2)$ and correlation $\rho$. The expectation $\mathbb{E}[\text{ReLU}(g_1)\text{ReLU}(g_2)]$ is given by the decomposition:
\begin{equation}
\label{eq:final_ntk_decomp}
\boxed{
\mathbb{E}[\sigma(g_1)\sigma(g_2)] = \sigma_1\sigma_2 \cdot I_{12} + \mu_1\sigma_2 \cdot I_{01} + \mu_2\sigma_1 \cdot I_{10} + \mu_1\mu_2 \cdot I_{00}
}
\end{equation}
where the terms $I_{ij}$ are moments of standardized bivariate normal variables $(Z_1, Z_2)$ with thresholds $a_1 = -\mu_1/\sigma_1$ and $a_2 = -\mu_2/\sigma_2$.
\begin{align*}
    I_{00} &= \mathbb{E}[\mathbf{1}_{Z_1>a_1, Z_2>a_2}] \\
    I_{10} &= \mathbb{E}[Z_1 \cdot \mathbf{1}_{Z_1>a_1, Z_2>a_2}] \\
    I_{01} &= \mathbb{E}[Z_2 \cdot \mathbf{1}_{Z_1>a_1, Z_2>a_2}] \\
    I_{12} &= \mathbb{E}[Z_1 Z_2 \cdot \mathbf{1}_{Z_1>a_1, Z_2>a_2}]
\end{align*}
For the symmetric case where $\mu_1=\mu_2=\mu$ and $\sigma_1=\sigma_2=\sigma$ (so $a_1=a_2=a=-\mu/\sigma$), the main formula simplifies to $\sigma^2 I_2 + 2\mu\sigma I_1 + \mu^2 I_0$, where $I_0=I_{00}$, $I_1=I_{10}=I_{01}$ and $I_2=I_{12}$. These terms have the analytical solutions:
\begin{itemize}
    \item \textbf{Term $I_0$ (Probability Term):} $I_0 = L(-a, -a; \rho)$
    \item \textbf{Term $I_1$ (First-Order Moment):} $I_1 = (1+\rho)\phi(a)\Phi\left(a\sqrt{\frac{1-\rho}{1+\rho}}\right)$
    \item \textbf{Term $I_2$ (Second-Order Cross-Moment):} $I_2 = \rho \cdot L(-a, -a; -\rho) + \phi(a, a; \rho)$
\end{itemize}
where the functions represent standard Gaussian integrals: $\phi(z) = \frac{1}{\sqrt{2\pi}}e^{-z^2/2}$ is the standard normal PDF, $\Phi(z) = \int_{-\infty}^z \phi(t)dt$ is its CDF, and $L(a, b; \rho) = \int_{-\infty}^a \int_{-\infty}^b \phi(z_1, z_2; \rho) dz_1 dz_2$ is the bivariate normal CDF.
\end{theorem}

\section{Future Work: The Multi-Layer Case}
The analysis for a deep MMNN with multiple layers builds upon the single-layer kernel derived here. It involves a recursive formula where the kernel at layer $L$ depends on the kernel from layer $L-1$. The detailed computation for this recursive case is left for future work.


\end{document}